%======================================================== 1 TITLE
\begin{frame}[plain]
\begin{tikzpicture}[remember picture,overlay]
\fill[soft] (current page.north east) rectangle ([xshift=-64mm]current page.south east);
\node[anchor=north west,align=left,text width=92mm] at ([xshift=9mm,yshift=-16mm]current page.north west) {%
  {\color{ink!55}\small Chuyên mục Thách thức Toán học · Đề ra kỳ này}\\[5mm]
  {\fontsize{22}{26}\selectfont\bfseries Hướng dẫn sử dụng}\\[3mm]
  {\color{ink!80}\normalsize Hệ thống quản lý bài đề xuất trên web}\\[7mm]
  {\footnotesize Phần chung cho mọi người, rồi một phần riêng cho từng vai trò:}\\[2mm]
  \rtag{TBT}\ \rtag{PT}\ \rtag{NCB}\ \rtag{PB}\ \rtag{VP}\ \rtag{BTK}\ \rtag{QT}\\[8mm]
  {\color{ink!55}\footnotesize Bản thử · kỳ 10/2026}};
% mini mock on the right
\begin{scope}[shift={([xshift=-60mm,yshift=26mm]current page.east)},x=1mm,y=1mm]
\browser{0}{0}{54}{50}
\node[anchor=west,font=\fsz{6}{7}\bfseries] at (2,-6.5) {Đề ra kỳ này};
\node[anchor=west,font=\fsz{3.6}{4.2},text=ink!60] at (2,-9.5) {ban-bien-tap@example.com · NCB};
\foreach \y/\c/\t/\m/\s/\k in {-13/tDS/ĐS/A/SL/2026-08-07a,-25/tHH/HH/B/Mới/2026-09-04b,-37/tSH/SH/A/SL-OK/2026-06-05a}{
  \draw[ink!20,fill=white,rounded corners=1pt] (2,\y) rectangle (52,\y-10);
  \fill[\c] (2,\y) rectangle (2.8,\y-10);
  \node[anchor=west,font=\fsz{4.4}{5}\ttfamily\bfseries,text=PB!70!black] at (4,\y-2.5) {\k};
  \node[pchip,anchor=west] at (22,\y-2.5) {\t}; \node[pchip,anchor=west] at (27,\y-2.5) {Mức \m}; \node[pchip,anchor=west] at (36,\y-2.5) {\s};
  \fill[ink!15] (4,\y-6) rectangle (48,\y-6.8); \fill[ink!15] (4,\y-8) rectangle (36,\y-8.8);}
\end{scope}
\end{tikzpicture}
\end{frame}

%======================================================== 2 MAP
\begin{frame}[label=map]{Hướng dẫn này gồm những gì}
\begin{tikzpicture}[x=1mm,y=1mm]
\draw[ink!25,fill=soft,rounded corners=3pt] (0,0) rectangle (66,-62);
\node[anchor=north west,font=\scriptsize\bfseries] at (3,-3) {PHẦN CHUNG — mọi người nên đọc};
\begin{scope}[every node/.style={anchor=north west,font=\fsz{6.6}{8}}]
\node at (3,-10) {\hyperlink{vao}{\textbf{1.}\ Vào hệ thống lần đầu}};
\node at (3,-16) {\hyperlink{taikhoan}{\textbf{2.}\ Đúng tài khoản Google}};
\node at (3,-22) {\hyperlink{duongdan}{\textbf{3.}\ Đường dẫn và phiên làm việc}};
\node at (3,-28) {\hyperlink{danhsach}{\textbf{4.}\ Danh sách bài và bộ lọc}};
\node at (3,-34) {\hyperlink{mabai}{\textbf{5.}\ Mã bài và thuộc tính}};
\node at (3,-40) {\hyperlink{trangbai}{\textbf{6.}\ Trang một bài}};
\node at (3,-46) {\hyperlink{nhanxet}{\textbf{7.}\ Viết nhận xét có công thức}};
\node at (3,-52) {\hyperlink{aithay}{\textbf{8.}\ Ai thấy gì}};
\end{scope}
\node[anchor=north west,font=\fsz{5.6}{6.8},text=ink!60,text width=60mm] at (3,-57.5) {Bấm vào một dòng để nhảy tới trang đó.};
\node[anchor=north west,font=\scriptsize\bfseries] at (72,-3) {PHẦN THEO VAI TRÒ — đọc phần của mình};
\foreach \r/\lab/\name/\what [count=\i] in {
  TBT/tbt/Tổng biên tập/{quyết định mức, tác giả, trạng thái; duyệt cuối},
  PT/pt/Phụ trách chuyên mục/{điều phối kỳ, xử lý xung đột, chọn bài},
  NCB/ncb/Người chuẩn bị bài/{soạn \LaTeX, đối chiếu bản gốc, ghi sửa đổi},
  PB/pb/Phản biện/{giải và nhận xét các bài được giao},
  VP/vp/Văn phòng/{nhận hồ sơ, giữ liên hệ tác giả},
  BTK/btk/BTK/{nhận bản \LaTeX\ khi khoá kỳ để chế bản},
  QT/qt/Quản trị/{thêm người dùng, giao vai trò}}{
  \node[anchor=west] at (73,{-12-(\i-1)*6.6}) {\hyperlink{\lab}{\rtag{\r}}};
  \node[anchor=west,font=\fsz{6.4}{7.6}\bfseries] at (85,{-12-(\i-1)*6.6+1.3}) {\hyperlink{\lab}{\name}};
  \node[anchor=west,font=\fsz{5.6}{6.6},text=ink!70] at (85,{-12-(\i-1)*6.6-1.6}) {\what};}
\node[anchor=north west,font=\fsz{5.6}{6.8},text=ink!60,text width=70mm] at (72,-57.5) {Một người có thể có nhiều vai trò; khi đó dùng được mọi quyền của các vai trò đó.};
\end{tikzpicture}
\end{frame}

%======================================================== 3 VÀO HỆ THỐNG
\begin{frame}[label=vao]{1.\ Vào hệ thống lần đầu}
\begin{tikzpicture}[x=1mm,y=1mm]
\foreach \i/\t in {1/{Mở đường dẫn đăng nhập},2/{Google xin phép biết email},3/{Kiểm tra email, bấm Vào hệ thống},4/{Danh sách bài của bạn}}{
  \node[badge=NCB,minimum size=5mm,font=\fsz{7}{8}\bfseries] at ({(\i-1)*37+3},-2) {\i};
  \node[anchor=west,font=\fsz{6.4}{7.4}\bfseries,text width=29mm] at ({(\i-1)*37+6.5},-2) {\t};}
% 1 email
\draw[win] (0,-8) rectangle (33,-40);
\node[anchor=north west,font=\fsz{5}{6},text width=30mm] at (1.5,-10) {\textbf{Email từ Pi}\\[1mm]Mời anh/chị phản biện kỳ 10/2026.\\Vào hệ thống tại:\\\textcolor{NCB}{\underline{script.google.com/…/exec}}};
\node[callout=ink,anchor=north west,text width=30mm] at (0,-43) {Đường dẫn này dùng cho mọi lần. Nên \textbf{đánh dấu trang} (bookmark).};
% 2 consent
\browser{37}{-8}{33}{32}
\node[font=\fsz{5}{6}\bfseries] at (53.5,-15) {Google};
\node[font=\fsz{4.6}{5.6},align=center,text width=29mm] at (53.5,-22) {„Pi – Đề ra kỳ này (đăng nhập)" wants to access your Google Account\\[1mm]\textcolor{ink!60}{• See your primary Google Account email address}};
\node[btn] at (60,-35) {Allow};\node[font=\fsz{4.6}{5.6},text=NCB] at (47,-35) {Cancel};
\node[callout=ink,anchor=north west,text width=30mm] at (37,-43) {Chỉ lần đầu. Ứng dụng chỉ được biết \textbf{địa chỉ email}, không gì khác. Bấm \textbf{Allow / Cho phép}.};
% 3 xin chào
\browser{74}{-8}{33}{32}
\fill[NCB!10] (74,-11.2) rectangle (107,-14);
\node[anchor=west,font=\fsz{3.1}{3.5},text=ink!70] at (74.6,-12.6) {This application was created by a Google Apps Script user};
\node[font=\fsz{3.8}{4}] at (105.3,-12.6) {\no};
\node[anchor=west,font=\fsz{5.4}{6.4}] at (76,-19) {Xin chào \textbf{ban@example.com}.};
\node[btn,anchor=west] at (76,-25) {Vào hệ thống};
\node[anchor=west,font=\fsz{4.2}{5},text=ink!60] at (76,-30.5) {Liên kết có hiệu lực 10 phút.};
\node[callout=warn,anchor=north west,text width=30mm] at (74,-43) {Nhìn email trong dòng \textbf{Xin chào}: đúng tài khoản của mình chưa? Dải xanh phía trên là thông báo của Google — bấm \no\ để đóng.};
% 4 list
\browser{111}{-8}{33}{32}
\node[anchor=west,font=\fsz{5.6}{6.6}\bfseries] at (112.5,-14) {Đề ra kỳ này};
\node[anchor=west,font=\fsz{3.6}{4.2},text=ink!60] at (112.5,-16.5) {ban@example.com · PB};
\foreach \y in {-19,-27}{\draw[ink!20,fill=white,rounded corners=1pt] (112.5,\y) rectangle (142.5,\y-6.5);
  \fill[ink!15] (114,\y-3.5) rectangle (139,\y-4.2);\node[anchor=west,font=\fsz{3.8}{4.4}\ttfamily\bfseries,text=PB!70!black] at (113.5,\y-1.6) {THU-01};}
\node[callout=ink,anchor=north west,text width=30mm] at (111,-43) {Dòng trên cùng: email và \textbf{vai trò} của bạn. Danh sách chỉ có những bài bạn được xem.};
\end{tikzpicture}
\vskip1mm
{\fsz{6.4}{7.6}\color{ink!75} Cần một tài khoản Google (Gmail). Không cần cài gì; dùng được trên máy tính và điện thoại. Nếu thấy „…chưa có vai trò trong hệ thống" — báo Phụ trách chuyên mục để được thêm vào.}
\end{frame}

%======================================================== 4 ĐÚNG TÀI KHOẢN
\begin{frame}[label=taikhoan]{2.\ Đúng tài khoản Google}
\begin{tikzpicture}[x=1mm,y=1mm]
\node[anchor=north west,text width=70mm,font=\fsz{7}{9}] at (0,0) {%
\textbf{Hệ thống nhận ra bạn qua tài khoản Google đang đăng nhập trong trình duyệt.}\\[2mm]
Nếu trình duyệt đăng nhập \textbf{nhiều tài khoản Google} (ví dụ Gmail cá nhân và Gmail công việc), Google có thể dùng
\textbf{tài khoản mặc định} — không phải tài khoản bạn định dùng.\\[3mm]
\textbf{Cách kiểm tra:} dòng \emph{Xin chào …} trước khi bấm \emph{Vào hệ thống}, và dòng email ở đầu danh sách bài.\\[3mm]
\textbf{Nếu sai tài khoản:}\\[1mm]
\hspace*{2mm}• mở \textbf{cửa sổ riêng tư} (Safari: File → New Private Window; Chrome: cửa sổ ẩn danh),\\
\hspace*{2mm}• chỉ đăng nhập đúng một tài khoản trong cửa sổ đó,\\
\hspace*{2mm}• mở lại đường dẫn đăng nhập.\\[3mm]
\textbf{Muốn đổi sang tài khoản khác lâu dài:} báo Phụ trách chuyên mục — vai trò gắn với địa chỉ email, không gắn với người.};
\begin{scope}[shift={(86,-2)}]
\browser{0}{0}{56}{38}
\node[anchor=west,font=\fsz{6}{7}] at (3,-10) {Xin chào \textbf{\textcolor{warn}{ca-nhan@example.com}}.};
\node[btn,anchor=west] at (3,-17) {Vào hệ thống};
\draw[warn,line width=1pt] (2,-7.5) rectangle (54,-12.5);
\node[callout=warn,anchor=north west,text width=50mm] at (3,-23) {Không phải tài khoản dùng cho Pi? \textbf{Đừng bấm} — đóng trang, mở cửa sổ riêng tư.};
\end{scope}
\node[callout=PB,anchor=north west,text width=54mm] at (86,-45) {Trên điện thoại: mở cùng đường dẫn trong Safari/Chrome; nếu điện thoại có nhiều tài khoản Google, kiểm tra dòng \emph{Xin chào} như trên.};
\end{tikzpicture}
\end{frame}

%======================================================== 5 ĐƯỜNG DẪN & PHIÊN
\begin{frame}[label=duongdan]{3.\ Đường dẫn và phiên làm việc}
\begin{tikzpicture}[x=1mm,y=1mm]
\node[hstep=NCB,text width=34mm,minimum height=12mm,anchor=north west,font=\fsz{6.4}{7.6}] (a) at (0,0) {\textbf{Đường dẫn đăng nhập}\\dùng mọi lần, đánh dấu trang được};
\node[hstep=PB,text width=34mm,minimum height=12mm,anchor=north west,font=\fsz{6.4}{7.6}] (b) at (54,0) {\textbf{Trang hệ thống}\\mở ra sau khi bấm Vào hệ thống};
\node[hstep=warn,text width=34mm,minimum height=12mm,anchor=north west,font=\fsz{6.4}{7.6}] (c) at (108,0) {\textbf{Phiên đã hết hạn}\\có nút \textbf{Đăng nhập lại}};
\draw[flow] (a) -- node[above,font=\fsz{5.4}{6.4}]{Vào hệ thống} (b);
\draw[flow] (b) -- node[above,font=\fsz{5.4}{6.4},align=center]{tải lại trang\\sau 10 phút} (c);
\draw[flow] (c.south) -- ++(0,-5) -| node[pos=0.25,below,font=\fsz{5.4}{6.4}]{Đăng nhập lại} (a.south);
\node[anchor=north west,text width=66mm,font=\fsz{6.6}{8.4}] at (0,-26) {%
\textbf{Nên}\\[1mm]
• Luôn vào từ \textbf{đường dẫn đăng nhập} (trong email, hoặc dấu trang).\\
• Cứ để tab mở mà làm việc: phiên dùng được \textbf{6 giờ}, miễn là không tải lại trang.\\
• Nếu gặp trang „hết hạn" — bấm \textbf{Đăng nhập lại}, mất vài giây.};
\node[anchor=north west,text width=70mm,font=\fsz{6.6}{8.4}] at (74,-26) {%
\textbf{Không nên}\\[1mm]
• Đánh dấu trang hay lưu địa chỉ của \emph{trang hệ thống} (dùng lại sẽ báo hết hạn).\\
• \textbf{Gửi địa chỉ đang hiện trên thanh địa chỉ cho người khác}: trong 10 phút đầu, địa chỉ đó cho phép vào hệ thống \emph{bằng tên của bạn}. Muốn mời ai, gửi \textbf{đường dẫn đăng nhập}.};
\end{tikzpicture}
\end{frame}
